\scanpage{9}

voit que $P(\mathcal M)\ne\varnothing\Longleftrightarrow\forall\alpha,\ P(M_\alpha)\ne\varnothing$. On est ainsi ramené au cas de \uline{$\otimes$-catégories semi-simples de type fini} \uline{sur} $\mathbb R$ : étudier l'ens. des polarisations d'une telle $\mathcal M$.

12) Soit $\mathcal M$ avec polarisation $P$ (sur $\mathbb R$ si on veut). Soit $M\in\operatorname{Ob}\mathcal M_{\mathbb C}^i$, i.e. un objet \inserted{$M_0$} de $\mathcal M^i$, avec action de $\mathbb C$ dessus\unsure{\,\textit{lui}|\,en} (i.e. endomorphisme $J_M$ de carré $-1$). Soit $\overline M$ avec l'action conjuguée ($J_{\overline M}=-J_M=J_M^{-1}$). Un hom

\[
M\otimes_{\mathbb C}\overline M\xrightarrow{\varphi}\mathbb C(-i)=\mathbb R(-i)\otimes_{\mathbb R}\mathbb C
\]

équivaut à un hom

\[
(M\otimes_{\mathbb C}\overline M)_{\underset{\substack{\inserted{\text{motif réel}}\\\text{sous-jacent}}}{0}}\longrightarrow\mathbb R(-i)
\]

i.e un hom

\[
M_0\otimes M_0\xrightarrow{\varphi_0}\mathbb R(-i)
\]

\[
\begin{aligned}
\text{satisfaisant}\quad\varphi_0(J_Mx,y)&=\varphi_0(\deleted{\unsure{X}}x,J_{\overline M}y)\\
&=-\varphi_0(x,J_My)\qquad\text{i.e.}\\
\varphi_0(J_Mx,J_My)&=\varphi_0(x,y)
\end{aligned}
\]

Pour que $\varphi$ soit \deleted{\unsure{soit}} ``hermitienne'' resp. ``antihermitienne'' il faut \unsure{s.|et s.|si} que $\varphi_0$ soit sym. resp. alternée [la symétrie hermitienne en $\varphi$ correspond à la symétrie ordinaire en $\varphi_0$]. On dit que $\varphi$ est une \uline{polarisation} \inserted{(pour $P$)} si $\varphi_0$ \deleted{est} est une (pour $P$). Cela implique $\varphi$ non dégénérée, i.e. $\widetilde\varphi:\check M\simeq\overline M(-i)$

% archive sheet number: 8
